% 第 3 题程序框图：按固定源 PDF 的框图结构逐一重绘（框形、分支标签与回边走向照原图）。
% 回边次序（易错点，600 dpi 逐箭头核对）：「是」分支先进入 \(s=s+k^{3}\)（箭头指向其下沿），
% 再上行进入 \(k=k+1\)，最后由 \(k=k+1\) 顶边折回主干，箭头指向主干（T 形汇入）。
% 坐标全部显式给定：主干 x=0；回边列 x=-3.4；join 在 init 底边下 0.35cm；
% inc 顶边(-2.6) 与 init 底边(-1.4) 留 1.2cm 空档，杜绝框体重叠与回边穿框。
\begin{tikzpicture}[>=stealth,every node/.style={font=\small},
  term/.style={draw,rounded corners,minimum width=1.5cm,minimum height=6mm},
  box/.style={draw,minimum width=2.2cm,minimum height=6mm},
  io/.style={draw,trapezium,trapezium left angle=70,trapezium right angle=110,minimum width=1.9cm,minimum height=6mm},
  decision/.style={draw,diamond,aspect=2.6,minimum width=2.6cm,minimum height=1.0cm,inner sep=0pt}]
\node[term] (start) at (0,0) {开始};
\node[box] (init) at (0,-1.1) {\(k=0,\ s=0\)};
\node[decision] (test) at (0,-6.6) {\(k\leqslant 2\)};
\node[io] (out) at (0,-8.0) {输出 \(s\)};
\node[term] (end) at (0,-9.2) {结束};
\node[box] (inc) at (-3.4,-2.9) {\(k=k+1\)};
\node[box] (add) at (-3.4,-4.5) {\(s=s+k^{3}\)};
\coordinate (join) at (0,-1.75);
\draw[->] (start)--(init);
\draw[->] (init)--(test);
\draw[->] (test.south)--node[right,xshift=2pt]{否}(out.north);
\draw[->] (out)--(end);
\draw[->] (test.west) -| node[above,pos=0.35]{是} (add.south);
\draw[->] (add.north)--(inc.south);
\draw[->] (inc.north) |- (join);
\end{tikzpicture}
